\section{讨论}
\label{sec:discussion}

\subsection{增益来自局部证据与视频覆盖}
Alpha-STALLED U0 是对 STALL 的局部化和视频级扩展，而不是可学习替代。冻结的 Global 分支保留整体外观与一阶转移真实性，Local 分支用细粒度同网格二阶 patch 轨迹补充局部动态异常。统一消融中，将 Local 加到 K=3 Global 后 Macro AP 提高 0.0237，将 Global 加到 K=3 Local 后提高 0.0431，说明两个分支都不可被另一个替代。K=3 相对统一单窗口的 Macro AP 提高 0.0087，视频级配对 bootstrap 区间为 [0.0057,0.0120]，说明额外时间覆盖带来稳定总体收益。

该收益并非所有数据集单调一致。K=3 在 ComGenVid 和 GenVideo 上分别提高 AP 0.0186 和 0.0119，但在 VideoFeedback 上下降 0.0045。这一结果限制了结论范围：三窗口主要修复长视频或间歇异常的单窗口遗漏，但对短视频占比高、开头窗口已充分或窗口间分布变化较大的集合，增加窗口可能稀释判别信号。

\subsection{为什么保持简单融合}
多项候选未达到预设的提升、最差数据集和生成器稳定性门槛。K=5、全部非重叠窗口、bottom-2 与 hybrid 聚合均未超过 K=3 分支均值；共同运动 mean/median residual 没有改善正式基线；14$\times$14 与 7$\times$7 token 的校准后融合不如 fine-only；中间层与最终层融合仅产生约 0.0002 AP 的变化；Mahalanobis 和 copula Joint Typicality 明显低于固定 $0.6/0.4$ 融合。上述结果说明，局部异常的稀疏性不足以抵消极值聚合噪声，共同运动也不是当前主要误差来源，而每数据集 200 个真实校准视频下的二维分布仍不足以稳定替代单调线性融合。

PatchSpatial 同样不是主要增益来源：在统一单窗口核心消融中，加入其 0.1 权重反而使 Macro AP 下降约 0.0070。U0 保留该项是为了与冻结 clean 定义一致，而非声称其独立有效；后续若以更严格的组件剪枝为目标，应优先验证删除 PatchSpatial，而不应继续扩大其权重。

\subsection{校准与协议边界}
U0 的全部白化参数、真实经验分布和冻结权重均不得使用生成视频或测试真实视频。每个校准视频在多窗口模式下只贡献一个视频级分数，并按 effective-$K$ 分组建模，避免长视频支配经验分布。校准规模实验显示，25/50/100/200 个真实视频对应的 Macro AP 均值为 0.8237/0.8487/0.8628/0.8680；200 个样本的跨 seed 标准差约 0.0022，且 Local 对校准规模比 Global 更敏感。因此，本文结论依赖目标真实域中约 200 个独立校准视频，不能直接外推到无目标域真实样本的部署。

跨数据集校准进一步验证了这一限制。Global 使用目标域和非目标域真实库时的平均 AP 分别为 0.8486 和 0.8475，几乎不变；Local 则从 0.8292 降到 0.7425。将三域 600 条真实视频合并后，最终 Macro AP 为 0.8636，仍低于目标域对角线的 0.8723。域依赖因此主要来自 Local patch 统计，而不是固定 $0.6/0.4$ 融合本身。更大的通用真实库可能缩小差距，但当前证据只支持目标域 real-only calibration。

锁定后的 GenVidBench 结果同时限定了“目标域校准”和“生成器泛化”的含义：使用 199 条新域真实视频校准后，U0 在此前未见的 ModelScope/Pika 评测上相对 STALL 提高 0.0460 AP，相对 clean K1 提高 0.0126 AP。该实验没有使用新域 fake 调参，因此支持统一结构在新生成器上的确认性迁移；但它仍使用新域真实视频，不能被解释为无目标域数据迁移。

历史 0.8737/0.8750 结果受校准泄漏影响，只作为协议审计记录，不参与方法选择或性能比较。历史 clean 单窗口 0.8570/0.8600 使用过旧的跨数据集选择过程；本文对 K=3 的因果增益以统一核心实现的 K=1 0.8632/0.8636 为直接对照。冻结主表 AP 以 real 为正类，而与原 STALL Table~1 对齐的扩展审计另报 generated-positive AP；二者不直接互换。配对重采样以视频 ID 为单位，避免窗口级伪重复。

生成器级平衡评测与数据覆盖回答不同问题。前者固定真假比例并让各生成器等权，但当生成视频数超过无泄漏真实池容量时，需要确定性选择 fake 子集；这不意味着校准集中的几百条真实视频就是完整测试集。duration-aware 协议实际保留 898/3,580/7,984 条测试真实视频，并保存全部 45,185 条生成视频分数。将全部 fake 纳入且用样本权重把 AP 类别先验固定为 50\% 后，Alpha 相对同协议 Global 的 Macro-3 AP$_{\rm fake}$/AP$_{\rm real}$ 仍提高 0.0368/0.0272，表明平衡子集不是观察到增益的必要条件。但固定先验只能消除类别比例对 AP 的机械影响，不能增加真实视频多样性。

相应地，1,000 次 cluster bootstrap 只量化固定测试身份上的有限样本不确定性，100 个选择 seed 检查进入平衡子集的具体身份是否重要，而全 fake 评测才检查已评分生成集合的覆盖率。三者均不能证明未收集视频或未来生成器上的性能。原 STALL 的生成器级 pairwise 协议同样会跨生成器复用真实视频；其 36,641 个平衡 pairwise real occurrences 来自约 6,822 个唯一测试 real，而不是 36,641 个互异真实文件。

原窗口因子审计进一步排除了“增益只来自换测试片段”的解释。全 45,185 条 fake 上，K1 STALL 到 K1 G+L 的 Macro-3 AP$_{\rm fake}$ 从 0.8188 提高到 0.8542，而继续改为 K3 G+L 后为 0.8629；在平衡身份 bootstrap 中，两步的平均增益分别为 0.0377 和 0.0085，95\% 区间均完全为正。由此可见，Local D2 是主要增益来源，三窗口覆盖是次要增益来源。三个 1 秒生成器的 K3 与 K1 完全相同，因此恢复它们不会人为放大多窗口收益；duration-aware 校准本身在 K1 的 AP$_{\rm fake}$ 差值区间跨 0，也不是总体改善的来源。

\subsection{局限性}
方法仍需要冻结 DINOv3 特征和目标域真实校准，计算量也会随有效窗口数近似线性增加。VideoFeedback 的小幅负迁移表明，均匀三窗口不是所有时长分布上的最优覆盖策略。锁定扰动实验中，CRF 23 和 10\% 丢帧的点损失较小且置信区间跨 0，但这不是等效性检验；CRF 35、25\% 重复帧、半分辨率恢复和 4 FPS 的 AP 差值区间均完全为负。同扰动 CDF 校准没有恢复严重退化，因此不能把部署漂移简化成分位数偏移。锁定注入实验进一步表明，Local D2 在自然生成视频上的统计互补性不等价于对任意编辑的单调异常响应：局部冻结只有很弱的正响应，而全帧冻结因二阶运动消失反而显著提高 Local 真实性。Patch-time map 对若干局部编辑的 AUPRC 也较低。因此 anomaly map 只能表示当前高斯典型性模型的分数贡献，不能被解释为通用语义伪影定位或因果解释。扰动结果必须与干净数据主结果分开理解。

当前正式比较覆盖 Original STALL、同核心 K1、两个单分支和锁定外部 GenVidBench，但尚未包含 D3、AEROBLADE、RIGID、ZED、T2VE、AIGVDet 等方法在相同视频交集、帧预算、校准数据和指标方向下的重跑。因此现有证据足以支持“相对 STALL 与同核心消融具有配对区间为正的改进”，不支持“优于所有现有方法”或“达到 state of the art”。若投稿目标要求竞争性 SOTA 声明，必须先补齐同协议外部基线；不同论文中的原始数字不能直接拼表替代。
